../cictikz/src/cictikz/data/tex/cictikz_preamble.tex

% Low frequency small signal model: transconductance g_m*v_gs and output
% resistance r_ds. Same frame as large_signal.tex, so the two figures read
% as before/after.

\begin{circuitikz}[american, thick, transform shape]

  ../cictikz/src/cictikz/data/tex/cictikz_lib.tex

  \def\ytop{2.6}
  \def\xsrc{2.2}
  \def\xrds{3.6}
  \def\xterm{5.0}

  % Gate terminal, floating
  \draw (0.09,\ytop) circle (0.09);
  \node[anchor=east] at (-0.1,\ytop) {G};
  \node[anchor=north] at (0.1,{\ytop-0.25}) {$+$};
  \node[anchor=west, color=blue] at (-0.15,1.3) {$v_{gs}$};
  \node[anchor=south] at (0.1,0.25) {$-$};

  % Transconductance
  \draw (\xsrc,\ytop) to [cI, l_=$g_m v_{gs}$] (\xsrc,0);

  % Output resistance
  \draw (\xrds,\ytop) to [R, l_=$r_{ds}$] (\xrds,0);

  % Rails to the terminals
  \draw (\xsrc,\ytop) -- (\xterm,\ytop);
  \fill (\xrds,\ytop) circle (0.075);
  \draw (\xterm+0.09,\ytop) circle (0.09);
  \node[anchor=west] at ({\xterm+0.3},\ytop) {D};

  \draw (0.0,0) -- (\xterm,0);
  \fill (\xrds,0) circle (0.075);
  \fill (\xsrc,0) circle (0.075);
  \draw (\xterm+0.09,0) circle (0.09);
  \node[anchor=west] at ({\xterm+0.3},0) {S};

\end{circuitikz}
\end{document}
